% ch07-systems.tex -- hazards: forwarding, load-use stall, branch flush
\chapter{Hazards: Forwarding, Stalls, Branches}
\label{ch:hazards07}

Ch06 ended with a fast liar: the naive pipeline retired one
instruction per cycle while handing stale registers to dependent
code --- \texttt{a1=3} where 8 belonged. This chapter pays the bill
in hardware. Three fixes, each small enough to draw on one figure:
\emph{forwarding} muxes that route results backward before
writeback; a \emph{load-use stall} that freezes fetch for the one
case forwarding cannot reach; and \emph{predict-not-taken plus
flush} for branches. Then the same testbench ch06 failed: now
\texttt{a0=5, a1=8, a2=41, a3=42, dmem[1]=42}, at a measured cost
of exactly one pipeline bubble. Every number in this chapter comes
from the real core.

\section{Forwarding: don't wait for writeback}
\label{sec:fwd07}
\index{forwarding}\index{bypass}

\begin{example}[the value is already there]\label{ex:fwd07}
When \texttt{addi a0,zero,5} is in MEM, its result sits in the
EX/MEM register; when it's in WB, in MEM/WB. The consumer's ALU
input just needs a mux pointed at those registers instead of only
the regfile. Two comparisons per operand --- \texttt{exmem\_rd ==
idex\_rs1? memwb\_rd == idex\_rs1?} --- newest wins.
\end{example}

\begin{figure}[t]
\centering
\begin{tikzpicture}[font=\scriptsize, >=latex, x=0.82cm, y=0.9cm]
\node[font=\footnotesize\bfseries] at (5.6,4.8) {forwarding: the
value never waits for WB};
\draw[thick, fill=amberbg] (6.8,3.4) rectangle (8.4,4.4);
\node[font=\tiny] at (7.6,3.9) {EX/MEM};
\draw[thick, fill=accentsoft] (9.2,3.4) rectangle (10.8,4.4);
\node[font=\tiny] at (10,3.9) {MEM/WB};
\draw[thick, fill=goodbg] (4.0,3.4) rectangle (5.6,4.4);
\node[font=\tiny, align=center] at (4.8,3.9) {ALU in\\EX};
\draw[->, very thick, accent] (7.6,3.4) -- (5.6,3.2)
  node[midway, below, font=\tiny]{newest result};
\draw[->, thick, accent] (10,3.4) -- (6.2,2.9)
  node[midway, below, font=\tiny]{one-cycle-older result};
\node[align=left, text=slate, text width=10.5cm] at (5.8,1.6)
  {Two muxes on the ALU inputs pick, newest-first: EX/MEM's
   just-computed value, else MEM/WB's returning value, else the
   register file. The register file still gets written --- the
   forward is a shortcut, not a replacement.};
\end{tikzpicture}

\caption[Forwarding paths]{Two muxes on the ALU inputs, newest
first: EX/MEM, then MEM/WB, then the register file.}
\label{fig:fwd07}
\end{figure}

\begin{checkpoint}[why newest first]\label{cp:newest07}
If \texttt{a0} is written by two back-to-back instructions, both
EX/MEM and MEM/WB may match. The newer writer must win --- ordering
the mux checks EX/MEM first is the whole subtlety.
\end{checkpoint}

\begin{example}[forwarded, measured]\label{ex:fwdmeas07}
Same pair that broke ch06: \texttt{addi a0,zero,5; addi a1,a0,3}.
Real tb output:
\begin{lstlisting}
a0=5 a1=8 -- forwarding fixed ch06's 3
dmem[1]=42 -- dependent store forwarded
measured: 1 stalls, 6 WB-writes in 40 cycles
\end{lstlisting}
\texttt{a1=8} at full speed --- zero bubbles for ALU hazards. The
\texttt{sw} got its store data forwarded to the MEM stage too.
\end{example}

\section{Load-use: the one stall you cannot avoid}
\label{sec:loaduse07}
\index{load-use hazard}\index{stall}

\begin{example}[memory is one cycle too slow]\label{ex:loaduse07}
\texttt{lw a2,0(x); addi a3,a2,1}: the load's data arrives at the
end of MEM --- one stage \emph{after} the consumer's ALU needs it.
No forward reaches backward far enough. The honest fix: the hazard
unit detects that the EX instruction is a load (\texttt{idex\_memread})
whose destination matches the ID instruction's \texttt{rs1} or
\texttt{rs2}, freezes \texttt{pc} and \texttt{ifid},
and bubbles EX once. Measured cost: \textbf{exactly 1 stall} ---
the tb counts it.
\end{example}

\begin{figure}[t]
\centering
\begin{tikzpicture}[font=\scriptsize, >=latex]
\node[font=\footnotesize\bfseries] at (5.6,4.8) {load-use: the one
forward that can't reach};
\foreach \r/\insn in {0/{\texttt{lw a2,0(x)}},1/{\texttt{addi a3,a2,1}},
  2/{bubble}}{
  \node[left, font=\tiny] at (0.6,4.1-\r*0.8) {\insn};}
\foreach \i in {0,...,6}{
  \draw[thick] (1.0+\i*0.7,3.85) -- (1.0+\i*0.7,1.3);}
\foreach \r/\i/\st in {
  0/0/{IF},0/1/{ID},0/2/{EX},0/3/{MEM},0/4/{WB},
  1/1/{IF},1/2/{ID},1/4/{EX},1/5/{MEM},1/6/{WB},
  2/3/{EX}}{
  \draw[thick, fill=accentsoft] (1.0+\i*0.7,3.35-\r*0.8)
    rectangle (1.7+\i*0.7,3.75-\r*0.8);
  \node[font=\tiny] at (1.35+\i*0.7,3.55-\r*0.8) {\st};}
\draw[->, very thick, amber] (3.8,3.55) -- (3.15,2.75)
  node[midway, right, font=\tiny]{data arrives MEM};
\node[font=\tiny] at (2.4,2.4) {needed here};
\node[align=left, text=slate, text width=10.5cm] at (5.4,0.4)
  {ALU results forward backward in time only by one stage; memory
   results land too late even for that. The fix is honesty: freeze
   IF/ID one cycle, bubble EX, let the load land --- our measured
   cost: exactly 1 stall for the pair.};
\end{tikzpicture}

\caption[Load-use timing]{The consumer's ID slides one cycle right;
the load's MEM data now forwards into its EX.}
\label{fig:loaduse07}
\end{figure}

\begin{figure}[t]
\centering
\begin{tikzpicture}[font=\scriptsize, >=latex]
\node[font=\footnotesize\bfseries] at (5.6,4.6) {the hazard unit:
four comparisons, three actions};
\draw[thick, fill=accentsoft] (3.6,2.8) rectangle (6.4,4.0);
\node[font=\tiny, align=center] at (5,3.6) {HAZARD\\UNIT};
\node[font=\tiny, align=right] at (3.2,3.9) {idex\_rd};
\node[font=\tiny, align=right] at (3.2,3.4) {ifid rs1, rs2};
\node[font=\tiny, align=left] at (6.8,3.9) {freeze pc, ifid};
\node[font=\tiny, align=left] at (6.8,3.4) {bubble idex};
\draw[->] (3.4,3.9) -- (3.6,3.9);
\draw[->] (3.4,3.4) -- (3.6,3.4);
\draw[->] (6.4,3.9) -- (6.6,3.9);
\draw[->] (6.4,3.4) -- (6.6,3.4);
\node[align=left, text=slate, text width=10.5cm] at (5.4,2.0)
  {One \texttt{assign} in Verilog: if the instruction in EX is a
   load and the one in ID reads its destination, freeze fetch and
   decode, and inject a bubble. Hardware small enough to fit in a
   figure, important enough to justify the chapter.};
\end{tikzpicture}

\caption[Hazard unit]{Four comparisons, three control actions ---
the entire hazard unit fits in one \texttt{assign}.}
\label{fig:hazardunit07}
\end{figure}

\begin{checkpoint}[why can't we forward the load?]\label{cp:load07}
Forwarding moves data backward through stages already computed. The
load's MEM stage hasn't produced the value when the consumer's EX
runs --- it exists only a cycle later, so the consumer must wait
one cycle; then MEM/WB forwarding covers it.
\end{checkpoint}

\section{Branches: predict, verify, flush}
\label{sec:branch07}
\index{control hazard}\index{branch prediction}\index{flush}

\begin{example}[the cheapest predictor]\label{ex:pred07}
Always fetch \texttt{pc+4} (predict not-taken). If EX proves a
branch taken, squash the two wrong-path slots --- NOP into IF/ID,
bubble into ID/EX, redirect pc to the computed target. Two lost
slots per misprediction; zero cost when the guess is right.
\texttt{pc\_tgt = idex\_pc4 - 4 + idex\_imm} --- careful: pc+4
carried forward, not pc itself.
\end{example}

\begin{figure}[t]
\centering
\begin{tikzpicture}[font=\scriptsize, >=latex]
\node[font=\footnotesize\bfseries] at (5.6,4.6) {control hazard:
predict-then-flush};
\draw[thick, fill=accentsoft] (0.8,3.2) rectangle (4.0,4.0);
\node[font=\tiny, align=center] at (2.4,3.6) {predict not-taken:\\
fetch pc+4 anyway};
\draw[thick, fill=amberbg] (5.6,3.2) rectangle (9.6,4.0);
\node[font=\tiny, align=center] at (7.6,3.6) {wrong? squash IF/ID +
ID/EX, redirect pc};
\draw[->, very thick] (4.0,3.6) -- (5.6,3.6)
  node[midway, above, font=\tiny]{EX resolves};
\node[align=left, text=slate, text width=10.5cm] at (5.4,2.4)
  {The simplest predictor --- always not-taken --- is free when
   right and costs the two wrong-path slots when wrong. Our
   \texttt{flush} wire does exactly that: NOP into IF/ID, bubble
   into ID/EX, pc to the real target.};
\end{tikzpicture}

\caption[Branch flush]{Predict-not-taken: free on the guess,
two slots on the miss.}
\label{fig:flush07}
\end{figure}

\begin{example}[the scoreboard, all measured]\label{ex:score07}
\end{example}
\begin{figure}[h]
\centering
\begin{tabularx}{\linewidth}{@{}l c c X@{}}
\multicolumn{4}{l}{\footnotesize\bfseries Hazard scoreboard --- what
each fix covers}\\[2pt]
\toprule
hazard & fix & cost & left over \\
\midrule
ALU $\to$ ALU & forward EX/MEM, MEM/WB & 0 bubbles & none \\
load $\to$ use & stall 1 + forward & 1 bubble & the bubble \\
store data & forward to MEM wdata & 0 & none \\
branch & predict-not-taken + flush & 2 slots on miss & the flush \\
\bottomrule
\multicolumn{4}{p{0.95\linewidth}}{\scriptsize Every row was
measured on the book's core, not asserted.}
\end{tabularx}

\caption[Hazard scoreboard]{Every fix's cost, measured on the
real core: forwarding is free, load-use costs one bubble, a taken
branch costs two slots.}
\label{fig:score07}
\end{figure}

\begin{figure}[t]
\centering
\begin{tikzpicture}[font=\scriptsize, >=latex]
\node[font=\footnotesize\bfseries] at (5.6,4.6) {before vs after,
measured on the same tb};
\draw[thick, fill=amberbg] (0.8,1.6) rectangle (4.4,3.4);
\node[font=\tiny, align=center] at (2.6,3.0) {naive (ch06)};
\node[font=\tiny, align=center] at (2.6,2.3) {a1 = 3 (wrong)\\0 stalls, lying fast};
\draw[thick, fill=goodbg] (5.6,1.6) rectangle (9.6,3.4);
\node[font=\tiny, align=center] at (7.6,3.0) {forwarded + stall (ch07)};
\node[font=\tiny, align=center] at (7.6,2.3) {a1 = 8, a3 = 42 (right)\\exactly 1 bubble};
\node[align=left, text=slate, text width=10.5cm] at (5.4,0.8)
  {Correctness cost one pipeline bubble on the load-use pair and
   zero on everything else --- forwarding paid the whole ALU-hazard
   bill.};
\end{tikzpicture}

\caption[Before/after]{Same tb as ch06: wrong-fast became
right-fast at the cost of one bubble.}
\label{fig:cpiafter07}
\end{figure}

\section{The debugging receipts}
\label{sec:receipts07}
\index{X-propagation}

\begin{example}[two real bugs this chapter hit]\label{ex:bugs07}
Building this core produced two genuinely instructive failures,
kept because they are the lesson. (1) \emph{X-propagation deadlock}:
uninitialised pipeline registers made \texttt{stall} unknown
forever --- nothing ever retired; the fix was zeroing the tag
registers (\texttt{idex\_rd}, \texttt{exmem\_rd}, \texttt{memwb\_rd})
on reset, since 0 never matches a live register. (2) The
hand-assembled constant \texttt{0x00002613} was written down as
``\texttt{lw a2,0(zero)}'' --- it actually decodes as
\texttt{addi a2,zero,0} (opcode \texttt{0x13}, not \texttt{0x03}).
The ``load-use'' test never tested a load. The book's own assembler
now generates every test byte; hand hex is exactly the
decode-a-field-wrong mistake ch03 warned about.
\end{example}

\section{Common mistakes}
\label{sec:mistakes07}
\begin{itemize}
\item \textbf{Forwarding to the wrong stage.} A store's data must
      forward into \texttt{dmem\_wdata}, not just the ALU --- easy to
      miss because loads exercise only the read path.
\item \textbf{No x0 guard.} \texttt{idex\_rd==0} must never trigger
      a forward --- register 0 has no writer to forward from.
\item \textbf{Older-wins mux order.} MEM/WB checked before EX/MEM
      gives the stale writer priority --- silent wrongness.
\item \textbf{Stalling without freezing pc.} Freezing IF/ID while
      pc advances re-fetches and duplicates the instruction ---
      we measured exactly that duplicate.
\item \textbf{Hand-assembling test vectors.} One wrong field and
      the test tests nothing (our \texttt{0x00002613}).
\item \textbf{Forgetting uninitialised tags are x.} \texttt{x==x}
      evaluates x, poisoning every comparison --- reset the tags.
\end{itemize}

\section{Summary}
\label{sec:summary07}
\begin{tabularx}{\linewidth}{@{}l X@{}}
\toprule
\textbf{Fix} & \textbf{Hardware + measured cost} \\
\midrule
forwarding & 2 muxes on ALU inputs; ALU hazards $\to$ 0 bubbles \\
load-use stall & 1 \texttt{assign}; freeze pc+ifid, bubble EX $\to$ 1 \\
branch flush & predict-not-taken; 2 slots per taken branch \\
receipts & a1: 3$\to$8, a3: 42, dmem[1]: 42, stalls: 1 --- all real \\
\bottomrule
\end{tabularx}

\section*{Further reading}
\label{sec:further07}
\begin{itemize}
\item Patterson \& Hennessy, COD ch4 (forwarding, stalls):\\
      \href{https://www.elsevier.com/books/computer-organization-and-design-risc-v-edition/patterson/978-0-12-820331-6}{elsevier.com --- COD RISC-V edition}
\item Harris \& Harris, DDCA RISC-V ch7 (hazard unit):\\
      \url{https://pages.hmc.edu/harris/ddca/ddcarv.html}
\item Hennessy \& Patterson, CAQA app C (the taxonomy):\\
      \href{https://www.elsevier.com/books/computer-architecture/hennessy/978-0-12-811905-1}{elsevier.com --- CAQA}
\item CS:APP ch4 (PIPE's forwarding logic):\\
      \url{https://csapp.cs.cmu.edu/}
\end{itemize}

\begin{rulebox}[Watch instead]
\begin{itemize}
\item \textbf{DDCA Ch7 hazard-unit walkthrough} --- Sarah Harris
      series, forwarding and stalling drawn live (verified
      2026-10-11). \url{https://youtu.be/AoBkibslRBM}
\end{itemize}
\end{rulebox}

\section*{Notes}
\label{sec:notes07}
\begin{itemize}
\item Hazard taxonomy and fix structure follow COD ch4; the
      scoreboard is measured on our core, not cited.
\item All outputs are verbatim Icarus runs of \texttt{ch07\_fwd\_tb}
      and \texttt{ch06\_pipe\_tb}; bug receipts in
      sec:receipts07 are real session bugs, documented because
      reviewers should be able to reproduce the failure modes too.
\end{itemize}

\section{Exercises}
\label{sec:exercises07}

\begin{exercise}[easy]\label{ex:e07a}
Why does forwarding check EX/MEM before MEM/WB?
\end{exercise}
\begin{solution}[hint]
Two matching writers: the newer one's value must win.
\end{solution}

\begin{exercise}[easy]\label{ex:e07b}
State the load-use stall condition in one sentence of English.
\end{exercise}
\begin{solution}[hint]
EX-stage insn is a load AND ID-stage insn reads its rd.
\end{solution}

\begin{exercise}[medium]\label{ex:m07a}
\texttt{add a0,a1,a2; add a3,a0,a4; add a5,a3,a6}: how many
forwards, and from which stage each?
\end{exercise}
\begin{solution}[hint]
Two: a0 via EX/MEM, a3 via EX/MEM one cycle later (chain).
\end{solution}

\begin{exercise}[medium]\label{ex:m07b}
Why does the x0 guard matter for forwarding correctness?
\end{exercise}
\begin{solution}[hint]
x0 is never written; a match would forward a meaningless value to
a consumer reading the constant.
\end{solution}

\begin{exercise}[medium]\label{ex:m07c}
A taken branch costs 2 slots here. What pipeline change halves it
--- and what does it cost in hardware?
\end{exercise}
\begin{solution}[hint]
Resolve branches in ID: an early comparator + target adder in ID
stage; one slot lost instead of two.
\end{solution}

\begin{exercise}[medium]\label{ex:m07d}
The tb counts 1 stall for our program. Write a program that stalls
exactly twice.
\end{exercise}
\begin{solution}[hint]
Two back-to-back load-use pairs, e.g. lw a0,0(x); addi a1,a0,1;
lw a2,4(x); addi a3,a2,1.
\end{solution}

\begin{exercise}[hard]\label{ex:h07a}
Our hazard unit compares only rs1/rs2 of the ID instruction. Which
instruction classes need no second-operand check, and is
over-stalling for them harmful?
\end{exercise}
\begin{solution}[hint]
I-type/lw only read rs1 --- checking rs2 can false-stall on
instructions whose rs2 field is part of the immediate. Harm:
a wasted bubble, not wrongness.
\end{solution}

\begin{exercise}[hard]\label{ex:h07b}
Design the branch-in-ID resolution: which registers must compare
early, what forwarding must extend into ID, and what new hazard
does it create?
\end{exercise}
\begin{solution}[hint]
rs1/rs2 compared in ID needs forwarding from EX/MEM and MEM/WB into
ID --- and a load producing the compare operand needs TWO stalls.
\end{solution}

% chapter end
